% Compact main-paper fragment. Detailed justification is in the supporting
% tensor_transfer_proof.tex. No standalone document or separate PDF is created.
\subsection{Product-law transfer without a Cartesian derivative maximum}
\label{sec:positive-tensor-transfer}

I complete the financial transfer for the four- and eight-asset models
by integrating derivative envelopes against a positive transport measure.
This retains dispersion moments instead of replacing $U$ by its largest
value on $[-L,L]^p$. Let $\nu$ be the unnormalised truncated Gaussian
measure, with mass $a$, and let $\tau=(a/N)\sum_c\delta_{c_c}$ use the
actual rounded nodes in the same CDF cells. For their centred cell moments
$\mu_{1,c},\mu_{2,c}$, Taylor's integral identity gives
\[
 \int f\,d(\nu-\tau)=\sum_c\mu_{1,c}f'(c_c)
                         +\int f''\,d\rho_2,
\]
where $\rho_2$ has nonnegative densities
$\int_{a_c}^{z}(z-y)\phi(y)dy$ to the left of $c_c$ and
$\int_z^{b_c}(y-z)\phi(y)dy$ to its right. Its mass is
$\sum_c\mu_{2,c}/2$. For
$\rho_1=\sum_c\widehat\mu_{1,c}\delta_{c_c}$ with certified
$\widehat\mu_{1,c}\geq|\mu_{1,c}|$, the absolute transfer is bounded
by $\int|f'|d\rho_1+\int|f''|d\rho_2$.

The exact moment-generating function of the second transport measure is
\begin{equation}
 \mathcal M_{\rho_2}(v)=\frac{
 e^{v^2/2}\{\Phi(L-v)-\Phi(-L-v)\}
 -(a/N)\sum_c e^{vc_c}-v\sum_c\mu_{1,c}e^{vc_c}}{v^2},
 \qquad v\ne0.
 \label{eq:positive-transport-mgf}
\end{equation}
This follows by substituting $f(y)=e^{vy}$ into the exact identity,
with continuous minus rounded measures in that order. At zero I use its
mass directly. Near zero, the enclosure
$e^{-|v|L}\rho_2(\mathbb R)\leq\mathcal M_{\rho_2}(v)
\leq e^{|v|L}\rho_2(\mathbb R)$ prevents cancellation in the quotient.

Define
$\Lambda_{j,k}=\nu^{\otimes(j-1)}\otimes\rho_k
\otimes\tau^{\otimes(p-j)}$ and
$\langle h\rangle_{j,k}=\int h\,d\Lambda_{j,k}$.
The product-measure telescoping identity then proves
\begin{equation}
 \left|\int R_2d\nu^{\otimes p}-\int R_2d\tau^{\otimes p}\right|
 \leq\sum_j\{\langle|\partial_jR_2|\rangle_{j,1}
                   +\langle|\partial_{jj}R_2|\rangle_{j,2}\}.
 \label{eq:positive-tensor-transfer}
\end{equation}
All raw masses remain inside the hybrid measures.

For $u=T_j$, $v=T_{jj}\geq0$ and $r=r_j$, let
$D_1\geq\|C_2'\|_\infty$, $D_2\geq\|C_2''\|_\infty$ and
$D_{31}\geq\|C_3'\|_\infty$, in addition to the shared $M_k$ bounds.
The integral representation of $R_2$ implies
\begin{align*}
 |\partial_jR_2|&\leq M_2U|u|+\tfrac12|r|D_1U^2
                                  +\tfrac16M_3U^2|u|,\\
 |\partial_{jj}R_2|&\leq M_2(u^2+Uv)+2|r|D_1U|u|
 +\tfrac23M_3Uu^2+\tfrac12r^2D_2U^2\\
 &\quad+\tfrac13|r|D_{31}U^2|u|
                +\tfrac16M_3U^2v+\tfrac1{12}M_4U^2u^2.
\end{align*}
To verify the coefficients, differentiate
$G=t^{-2}C_2(m+\log t)$ with $t=1+zU$:
\[
 G_j=rt^{-2}C_2'+zu t^{-3}C_3,
\quad G_{jj}=r^2t^{-2}C_2''+2rzu t^{-3}C_3'
                         +zv t^{-3}C_3+z^2u^2t^{-4}C_4.
\]
Apply the product rule to $U^2G$, use $t\geq1$, and integrate
$1-z$, $z(1-z)$ and $z^2(1-z)$, whose integrals are $1/2$, $1/6$
and $1/12$. Cauchy's inequality gives
$\langle U|u|\rangle\leq\sqrt{\langle U^2\rangle\langle u^2\rangle}$
and
\[
 \langle U^2|u|\rangle\leq
 \min\{\sqrt{\langle U^2\rangle\langle U^2u^2\rangle},
          \sqrt{\langle U^4\rangle\langle u^2\rangle}\}.
\]
These inequalities hold for finite positive measures without a mass
divisor. Every remaining moment is a polynomial of degree at most four
in the asset exponentials. Its signed multinomial sum factorises into
scalar MGFs, using $\rho_k$ in coordinate $j$, $\nu$ before it and
$\tau$ after it. Thus no $N^p$ residual values enter the certificate.
The outside-box $U^2$ moment follows by Gaussian tilting and contributes
$M_2\mathbb E[U^2\mathbf1_{Y\notin[-L,L]^p}]/2$.

For a fair first-correction comparison, I also transfer
$R_1=R_2+UP'(m)$. The extra derivatives are
$uP'+rUP''$ and $vP'+2ruP''+r^2UP'''$. Capped kernels use the shared
global Greek bounds; vanilla kernels use
$P'\leq Ke^{m+\lambda^2/2}$, $|P''|\leq Ke^{m+\lambda^2/2}+M_2$
and $|P'''|\leq Ke^{m+\lambda^2/2}+3M_2+M_3$.
The corresponding tilted positive-measure moments also factorise.
Its tail is bounded by the tilted outside moment of
$Ke^{m+\lambda^2/2}U$, with the additional $H$ tail bound for a cap.
I replace the second-residual tail by this first-residual tail.

At unchanged $L=6$, $b_y=10$, $q=28$, all 24 four- and eight-asset
second-correction transfer bounds pass $0.005$; their largest allowance
is $0.002284$ rounded upwards. All 24 first-correction competitors
also pass; their largest allowance is $0.004914$ rounded upwards.
Each global-envelope family retains three two-asset allocation failures,
for which the earlier sharper local certificate remains available.
The transport argument applies to arbitrary finite positive basket
weights and bounded boundary derivatives, beyond these frozen models.
Its numerical scope and polynomial moment-construction cost do not
establish a quantum advantage or a literature-priority claim.
